% GDZ PPN236018647,native576upper;HKUST577.
\clearpage\SourceLeaf{576}\ArithmeticHead{547}{R}{1800 MAI. 6—1800 MAI. 22.}
\NumberedItem{106}
\indent Incrementum ingens huius theoriae Brunov. Mai. 22 invenire contigit, per quod simul omnia praecedentia nec non theoria mediorum arithmetico-geometricorum pulcherrime nectuntur infinitiesque augentur.
\par\noindent\hfill[Brunovici, 1800 Mai. 22.]

\medskip\begingroup\small
\indent Es handelt sich hier um die in der Scheda Ac, S. 26 ff. aufgezeichneten Entdeckungen, die oben im Abschnitt [V.] der \textit{Theorie des agM.}, S. 194 ff. artt. [1.]—[8.] abgedruckt sind. »Omnia praecedentia« bedeutet die Theorie der lemniskatischen Funktionen. Der Zusammenhang mit dem agM. tritt besonders im art. [5.] hervor; er besteht darin, daß die im art. [3.] aufgestellten Formeln für die Transformation zweiter Ordnung der Funktionen
\[
T\varphi,\ W\varphi,\ T\left(\frac\varpi2-\varphi\right),\ W\left(\frac\varpi2-\varphi\right)
\]
\noindent (die im wesentlichen nichts anderes sind als die \textsc{Jacobi}schen Thetafunktionen), wenn man \(\varphi=0\) setzt, für die Reihen
\[
W0,\ T\frac\varpi2,\ W\frac\varpi2
\]
\noindent die Beziehungen des agM. liefern. Der Zusammenhang zwischen dem agM. und den Reihen, deren Exponenten die Quadratzahlen sind, erscheint also hier gleichsam als die Projektion der Transformationsformeln des art. [3.] auf die Ebene \(\varphi=0\).
\par\noindent\hfill\textsc{Klein.\quad Schlesinger.}
\par\endgroup
